\section{Introduction}
\label{sec:intro}

LLM applications routinely place text they do not control into the model's input: retrieved web pages, emails,
tables, code, and the outputs of tools called by an agent. Instructions hidden in that text can redirect the model,
a threat known as indirect prompt injection (IPI)
\citep{greshake2023signed,liu2023formalizing,yi2023bipia,zhan-etal-2024-injecagent,debenedetti2024agentdojo}.
A common line of defense is a detector: a classifier that reads the external text, possibly together with the
user's request, and decides whether it contains an injection before it reaches the model
\citep{li-etal-2025-piguard,ivry2025sentinel,chennabasappa2025llamafirewall,liu2025datasentinel}.

This paper is about what such detectors are asked to learn, and what their training and evaluation data let them
learn instead. Two observations motivate it.

\paragraph{Injection is relative to the request.}
Consider the sentence \emph{``Write the reply so it confirms the seat by Friday''} inside an email. If the user asked
the assistant to draft a reply accepting the invitation, following it serves the user. If the user asked for a
summary of the email, following it replaces the summary with a reply. The text is identical; only the request
differs. BIPIA defines indirect injections by the deviation they cause from what the user expects
\citep{yi2023bipia}, which makes the label a function of the request as well as of the content. Yet most deployed
detectors read the content alone, and the benchmarks used to train and compare them label every inserted attack
string as an injection whatever the user asked. A detector that ignores the request cannot separate the two cases
above, and existing evaluations cannot tell whether it does.

\paragraph{Insertion-built data carry construction artifacts.}
IPI training and test sets are usually built by inserting attack strings, and benign look-alikes, into clean
documents. Every choice in that procedure (where the string goes, how it is joined to the document, which templates
supply the negatives) can correlate with the label. Prior work has shown that IPI evaluations admit content-free
shortcuts in agent conversations \citep{kale2026probes} and that defenses learn position and trigger-token
heuristics \citep{li2026surface}. We find the same failure in BIPIA-style insertion data, measure it, and remove it.

\paragraph{Contributions.}
\begin{enumerate}
  \item \textbf{Request-counterfactual intent pairs} (\S\ref{sec:pairs}). Each pair holds a document and one inserted
  instruction byte-identical and varies only the user request: under one request following the instruction serves
  the user (label 0), under the other it subverts the request (label 1). Pairs are built as crossed quads, so that
  both a content-only detector and a request-only model score exactly chance by construction. The pairs are
  LLM-drafted, and only pairs that two independent human annotators label as intended are used.
  \todo{annotation not yet done; second annotator not yet identified}
  \item \textbf{Intent-Counterfactual Training (\ict{})} (\S\ref{sec:ict}). A text-only query/context cross-encoder
  trained on artifact-controlled data plus counterfactual pairs, with a pairwise margin loss between the two members
  of each pair. It needs no white-box access to any LLM. An ablation with the same backbone and the request as input,
  but without pairs or margin loss, separates the effect of the training signal from the effect of seeing the request.
  \item \textbf{An inventory of construction artifacts in BIPIA-style insertion data} (\S\ref{sec:artifacts}).
  Content-free probes (the insertion position alone, an untrained first-paragraph rule, and nine surface features of
  the inserted span) predict the label well above chance, and position-controlled stress tests show that fine-tuned
  detectors rely on insertion position. This extends \citet{kale2026probes} from agent conversations to
  insertion-based document data.
  \item \textbf{An artifact-controlled dataset (\vthree) and a generalization study} (\S\ref{sec:v3}). \vthree{}
  samples position and separator independently of the label, preserves document formatting, matches the surface form
  of benign and injected spans, holds out attack categories for model selection and attack files for testing, and
  splits by near-duplicate document component. It passes pre-registered content-free gates. We ask whether training
  on it, rather than on artifact-laden data, improves transfer to agent benchmarks, extending \citet{liu2026passing}
  to training on BIPIA-style data.
\end{enumerate}

\paragraph{Pre-registration.}
Our four hypotheses (H1: artifact removal improves transfer; H2: \ict{} reads the request; H3: \ict{} beats
request-aware comparators on intent-dependent sets and is non-inferior on content-only transfer sets; H4: \ict{}
raises fewer false alarms on requested instructions), their success criteria, the evaluation sets, the threshold
protocol and a decision table that fixes which kind of paper results lead to were registered before any \vthree{},
pair, or agent-benchmark result existed. Every pre-registered comparison is reported, whatever its outcome, and every
post-hoc change is listed as a deviation. \todo{cite the registration (commit hash / OSF link) and summarize which
row of the decision table the results fall in}
